\section{Otimização da Regressão}

Essa formulação já basta para viabilizar a modelagem computacional.
Ainda assim, convém enxergar onde estão os gargalos antes de escrever uma linha
de código.
O sistema linear proposto tem alta complexidade, e vale a pena organizar o
raciocínio desde já.
Seja $N_r$ o número de termos $\hat{e}_d\,\vec{x}^{\vec{\alpha}}$ de grau
igual a $r$ ($\|\vec{\alpha}\| = r$); então
\begin{equation*}
    N_j
    \; = \;
    m \cdot \#\{\vec{\alpha}\in\mathbb{N}^n:\ |\vec{\alpha}|=j\}
    \; = \;
    m\,\binom{n+j-1}{j}.
\end{equation*}

Assim, podemos calcular o número de termos existentes de grau até $r$
\begin{equation*}
    \| \{ \varphi_k \}\|
    \;=\;
    N_{\le r}
    \;=\;
    \sum_{j=0}^{r} N_j
    \;=\;
    m \sum_{j=0}^{r} \binom{n+j-1}{j}
    \;=\;
    m\,\binom{n+r}{r}.
\end{equation*}

Ou seja, a nossa matriz $G$ tem dimensão considerável:

\begin{equation*}
    N_{\le r} \times N_{\le r}
    = \bigl(m\,\binom{n+r}{r}\bigr)^2
    = m^2 \binom{n+r}{r}^2
    = m^2 \binom{n+r}{n}^2
    = m^2\,\frac{(n+r)!^{\,2}}{(n!)^{2}(r!)^{2}}.
\end{equation*}

Isso nos dá a complexidade para a regressão polinomial:

\[
    p=\binom{n+r}{r}.
\]

Via método QR:
\[
    T_{\mathrm{QR}}(N,n,r,m)=O\!\left(Np^{2}+p^{2}m\right)
\]

Via equações normais (Cholesky):
\[
    T_{\mathrm{NE}}(N,n,r,m)=O\!\left(Np^{2}+p^{3}+p^{2}m\right)
\]

Essas estimativas refletem a comparação clássica entre decomposição QR e
equações normais/Cholesky na resolução de mínimos
quadrados\cite{burden_faires_2011,franco_calculo_numerico_2006}.

Os fatoriais elevados a 2 ou 3 já são razão suficiente para deixar qualquer
cientista da computação em alerta.
Não conseguimos melhorar muito a complexidade assintótica desse algoritmo,
porém é possível extrair um ganho de desempenho prático explorando uma
propriedade estrutural de $G$.
Note que, por definição,
$$
\hat e_i \perp \hat e_j \quad i \ne j \quad \implies
\langle \hat e_i \; c_{d,\vec{\alpha}} \; \vec x ^{\vec \alpha},
\hat e_j \; c_{d,\vec{\beta}} \; \vec x^{\vec \beta}\rangle = 0
$$

Então, com uma boa escolha de mapeamento do índice $k$, nossa matriz $G$
torna-se diagonal por blocos:

$$
A \;=\;
\begin{array}{ccc}
  &

    e_1 \quad\;e_2 \quad\; \cdots \quad\; e_n

  & \\[1ex]
%
  \left. \,
    \begin{matrix}
      e_1\\
      e_2\\
      \vdots\\
      e_h
    \end{matrix}
  \right.
  &
    \begin{bmatrix}
      A_1      & 0        & \cdots & 0      \\
      0        & A_2      & \cdots & 0      \\
      \vdots   & \vdots   & \ddots & \vdots \\
      0        & 0        & \cdots & A_h
    \end{bmatrix}
  &
  \left.
    \begin{matrix}
      \\ \\ \\
    \end{matrix}
  \right.
\end{array},
$$

Usando o delta de Kronecker, podemos resumir essa propriedade da seguinte forma:

$$
    G_{ij} =
    \langle \varphi_i ,\, \varphi_j \rangle \cdot
    \delta_{i ,\, j}
    ,\quad\quad
    \hat e_i \perp \hat e_j
$$

Logo, em vez de resolver um sistema linear de tamanho catastrófico, exploramos
as propriedades de matrizes diagonais por bloco e decompomos o problema em $m$
sistemas menores do tipo
$$A_i \, a_i = b_i.$$

Essa estratégia não altera a complexidade assintótica, mas proporciona um ganho
de desempenho prático relevante.
Além disso, como as matrizes $A_i$ são idênticas entre si --- variando apenas o
vetor unitário associado a cada componente ---, é possível reutilizar
artefatos intermediários e paralelizar os solvers sem esforço adicional.

